% Prefer a top-of-column slot so the figure can precede the next section.
\begin{figure}[!t]
\centering
% Vector drawing generated by tools/figures/build.py; rerun it after changing the figure.
\ifbilingual
\includegraphics[width=\linewidth]{figures/two-schemas-zh}
\else
\includegraphics[width=\linewidth]{figures/two-schemas}
\fi
\caption{\bt{End-to-end accuracy on two schemas, \llmmodel{}. Groups pool the held-out questions, four non-breaking updates, five breaking changes covered by the conditions, and all twelve settings, with equal weight per setting. Bars show pooled accuracy; error bars span the per-run accuracies (five runs for synthetic retail, \GeRuns\ for TPC-H SF1).}{两种模式上的端到端正确率，\llmmodel{}。各组分别合并留出题、4 种正常更新、条件覆盖的 5 种破坏性变化和全部 12 种情形，每种情形等权。柱高为合并正确率；误差线表示各次运行正确率的最小值到最大值（合成零售数据为 5 次，TPC-H SF1 为 \GeRuns\ 次）。}}
\label{fig:twoschemas}
\Description{Controlled comparison of no memory and MAVRA on two schemas, with accuracy from zero to 100 percent. Synthetic retail, five runs: held-out 58 and 100 percent; non-breaking updates 49 and 98; covered breaking changes 32 and 72; all settings 41 and 80. TPC-H SF1, three runs: held-out 84 and 96; non-breaking updates 83 and 97; covered breaking changes 60 and 75; all settings 71 and 83. Capped error bars span the per-run accuracies. Published memory methods are compared in the separate RelBench study.}
\end{figure}
